% RESUMEN

\VIUSectionStar{Resumen}

Consistente en una breve descripción (menos de una página, generalmente de 200 a 300 palabras) del trabajo a realizar, con sus requisitos y especificaciones, mencionando explícitamente si se basa en trabajos previos realizados por el tutor del alumno, un proyecto anterior o similar.

Las faltas de ortografía no son admisibles en un trabajo académico como un TFM asociado a Estudios Superiores.

Un TFM no es un diario personal. Trata de evitar la primera persona. Por ejemplo, en lugar de, “En mi opinión el modelo xxxxxx es muy bueno...” cambiar a “A la vista de los datos analizados o fuentes consultadas el modelo está obteniendo buenos resultados…”. Está permitido dar tu opinión siempre que lo precise y esté argumentada.

La claridad en la redacción del texto para la compresión del mismo es fundamental. Trata descomponer oraciones excesivamente compuestas en simples. Evita en la medida de lo posible oraciones subordinadas.

Evita usar estructuras generales o estereotipos del tipo “Como afirman numerosos expertos…” y sustitúyelos por “De acuerdo a los datos o información obtenidos …”

Evita usar las expresiones demasiado coloquiales, el humor o la ironía.

Usa sinónimos para evitar repetir el mismo término varias veces. Lee continuamente el texto para evitar contenido repetido o inconsistente.

Se recomienda incluir una versión en inglés(Abstract) al final del resumen.

\textbf{Palabras clave:} Entre 4 y 8 palabras clave(keywords).
